\section{Introduction}

Finite-context world models often fail in a way that is easy to miss: they can
observe a decisive cue early and later produce a representation that behaves as
if the cue was never present.  This matters for physical intelligence because
action-relevant evidence is rarely synchronized with the moment of action.  A
robot can read a sign before turning away, see a door state before entering a
hallway, or observe an object identity before that object becomes occluded.  If
the interface to a frozen host exposes only a recent window, then locally
plausible prediction is not enough; the old evidence must survive long enough,
in a usable form, to change a later decision.

We study this failure as an \emph{interface} problem rather than a raw memory
capacity problem.  The question is not only whether a hidden state can memorize
something.  The stricter question is whether old visual evidence, after it has
left the legal context window, is (i)~\emph{retained} in a causal memory state,
(ii)~\emph{exposed} through a host-facing representation, and (iii)~\emph{usable}
by a bounded downstream controller.  These three claims are routinely conflated.
A model may store old evidence internally yet fail to expose it in its output.
A linear readout may decode a label through a shortcut if the current endpoint
still leaks the cue.  A native controller may succeed only because it received
the true goal rather than the memory-selected one.  Our evaluation is designed
to close each of these loopholes with matched controls.

We introduce a \emph{masked-evidence slot-memory interface}
(Figure~\ref{fig:architecture}).  Training is self-supervised: it never receives
cue labels, cue crops, or manually selected key frames.  A target miner divides
the causal render stream into temporal bins, scores visual saliency, selects
salient patch sets, and removes them from the visible stream.  A compact slot
transformer must then predict the stop-gradient latent targets of those hidden
patches from a small persistent memory state.  Cue labels enter only
\emph{after} training, when a post-hoc ridge readout tests whether the learned
memory exposes the old cue under matched full, reset, and no-state conditions.

The central empirical point is structural rather than merely numerical.  Under
the same caches, target mining, patch target encoder, loss, optimizer, and
readout, replacing the slot writer with GRU, LSTM, or Mamba-lite recurrence
weakens long-delay exposure.  The recurrent baselines often solve short,
waypoint-like rows, which is useful evidence that the task is not impossible;
their failure appears when the cue is old or distributed across object-like
visual evidence.  Slots help because the target itself is a \emph{set} of visual
facts, and a memory set is a better carrier for that target than a single
recurrent vector.

Concretely, on twelve OGBench render decks, three controlled evidence ages
(4, 8, 15), and three seeds, the slot writer passes every registered retention
cell (36/36 environment-age rows, 108/108 seed cells) with mean full-memory
balanced accuracy $0.983$ while reset and no-state controls stay at $0.250$ and
$0.248$.  The strongest recurrent baseline reaches $0.803$ mean full accuracy
and $22/36$ rows; at age~15 the slot writer scores $0.962$ against $0.669$,
$0.350$, and $0.683$.  Where audited fixed controllers exist, memory-selected
targets execute at $0.957$ mean success across $27/36$ rows.

\paragraph{Contributions.}
\begin{itemize}
  \setlength\itemsep{1pt}
  \item A \textbf{fail-closed audit} for old-evidence memory in frozen
    world-model interfaces, separating retention, host exposure, and bounded
    use with matched reset / no-state controls.
  \item A \textbf{label-free masked patch-set JEPA objective} that trains a slot
    memory writer from automatically mined salient patches, with full pseudocode
    (Algorithms~\ref{alg:train}--\ref{alg:audit}).
  \item A \textbf{controlled architecture comparison} across twelve OGBench
    decks, three evidence ages, and three seeds, isolating the writer format
    while holding every other component fixed, extended with capacity-matched
    and structured-memory baselines (SlotSSM-style, gated slot attention,
    top-$k$ retrieval, cached attention).
  \item A \textbf{bounded streaming formulation} $M_t=f_\theta(M_{t-1},
    Z_{t-K+1:t})$ with token eviction and constant per-step memory/compute, plus
    a \textbf{delay-generalization} and \textbf{age-scaling} study to age~128.
  \item A \textbf{frozen-LeWM integration} showing the same interface exposes
    old evidence in the host's own output beyond its context window, and a
    \textbf{shortcut-broken benchmark} confirming retention of object identity
    rather than colour.
\end{itemize}

\noindent
Our contribution is best read as a \emph{fail-closed evaluation protocol plus a
controlled demonstration that target--carrier structural alignment improves
host-visible old-evidence exposure}, rather than a wholly new slot-memory
architecture: slot memory \citep{r_slotssm,r_gsa,r_slotmemory}, masked JEPA
prediction \citep{r_causaljepa}, and memory-augmented world models
\citep{r_onmemory,r_lewm} each exist; the novelty is the interface audit that
separates storage, host exposure, and use and shows where structural alignment
matters.

\section{Related Work}

\paragraph{World models under memory pressure.}
Latent world models compress observations into predictive state for control
\citep{r_dreamer}.  Recall-to-Imagine makes long-term memory and credit
assignment explicit bottlenecks for model-based RL \citep{r_r2i}.  Diffusion and
video world models face a related pressure: short observation windows can
produce plausible futures that drift from earlier evidence.  LoopNav turns this
into looped spatial-consistency tests \citep{r_loopnav}; StateSpaceDiffuser,
long-context state-space video models, and memory-expert diffusion world models
add state-space or expert memory paths to preserve temporal consistency across
long rollouts \citep{r_statespacediffuser,r_lcsvwm,r_memory_experts}.  We ask a
narrower but stricter interface question: not whether a generator looks
temporally consistent, but whether old evidence that is illegal for the current
context is still exposed in a frozen host-facing representation.

\paragraph{Partial observability and memory in RL.}
Partially observable RL has long treated memory as a substitute for unavailable
Markov state.  PO-Dreamer adds memory guidance to world models under partial
observability \citep{r_po_dreamer}, and memory traces replace fixed observation
windows with compact, analyzable history summaries \citep{r_memory_traces}.
These works optimize agents or value estimates.  In contrast, our audit isolates
a representation-interface claim: the downstream readout receives the same
endpoint under full, reset, and no-state conditions, so success must come from
retained old evidence rather than a policy learning to exploit a different
observation stream.

\paragraph{Sequence memory and learned state.}
GRU and LSTM recurrence remain standard baselines for compressing sequence
history \citep{r_gru,r_lstm}.  Mamba and Mamba-2 scale selective state-space
recurrence and connect it to attention-like computation
\citep{r_mamba,r_mamba2}.  Infini-attention, test-time-training layers, and
Titans push the fixed-context problem from compressive attention, self-supervised
hidden-state updates, and neural long-term memory
\citep{r_infini_attention,r_ttt,r_titans}.  SlotSSM and Gated Slot Attention are
closest to our writer because they make sequence memory modular or slot-like
rather than a single monolithic vector \citep{r_slotssm,r_gsa}.  These works
support the broad argument that fixed-window attention is insufficient for
long-horizon evidence; they do not, by themselves, answer whether a memory state
exposes the task-relevant old cue under leakage controls, which is exactly what
we hold fixed and test.

\paragraph{Structured and embodied memory.}
Object-centric world models show that decision-relevant visual details are often
small, structured, and easily washed out by reconstruction losses
\citep{r_oc_storm}.  SOLD and language-guided object-centric world models show
that slot latents can serve as relational state interfaces for manipulation
\citep{r_sold,r_object_wm_lgm}.  Embodied-memory systems such as KARMA,
Retrieval-Augmented Embodied Agents, Memo, and 3DLLM-Mem separate short-term
working context from longer experience or spatial-temporal memory
\citep{r_karma,r_raea,r_memo,r_3dllm_mem}; agent-memory work such as Memory-R1
and ReMemR1 studies what to store, update, and retrieve
\citep{r_memory_r1,r_rememr1}.  Evaluation work such as RULER and LongMemEval
warns that easy retrieval tests overstate real memory ability
\citep{r_ruler,r_longmemeval}.  Our slot memory is aligned with this structured
direction, but the supervision and claim differ: we train from label-free masked
visual patch sets and reveal labels only to audit whether stored evidence is
linearly accessible and physically useful.

\section{Problem Formulation}

Let $x_{0:T}$ denote rendered frames and $a_{0:T-1}$ actions.  A cue $c \in
\{0,1,2,3\}$ is visible at an early time and determines a later target label
$y=c$.  At readout time $t$, the frozen host sees only a recent context window
of length $K$; the cue lies strictly outside that window.  We call the delay
between the cue's last visible frame and the readout boundary the
\emph{evidence age} $\alpha$.  The registered ages are $\alpha\in\{4,8,15\}$,
acting as controlled stress tests rather than automatically discovered horizons.

\paragraph{Audit conditions.}
The audit has three conditions that share the same held-out episodes, endpoint,
and readout, differing only in the memory state fed to the readout.
The \emph{full} condition uses the causal memory state $M_t$ produced from the
full legal stream.  The \emph{reset} condition observes the same endpoint and
recent context but clears persistent memory at the context boundary.  The
\emph{no-state} condition removes persistent memory entirely.  Because all three
share the endpoint, any information recoverable from current pixels, the action
trace, class imbalance, or a forgiving readout raises \emph{all} conditions
together and therefore fails the audit.

\paragraph{Claim ladder.}
We report a four-rung ladder in which each rung is interpreted only after the
previous one passes.  \textbf{(1)~Demand:} the decisive cue must sit outside the
legal context window (ages $4/8/15$).  \textbf{(2)~Retention:} full memory must
expose the old cue on held-out episodes.  \textbf{(3)~Controls:} reset and
no-state features must stay near chance.  \textbf{(4)~Use:} executed success is
reported only where a fixed local controller exists; missing rows are marked
controller-unavailable rather than counted as successes.  Formally, a seed cell
passes when
\begin{equation}
\mathrm{BAcc}_{\text{full}} \ge 0.750
\quad\wedge\quad
\mathrm{BAcc}_{\text{reset}} \le 0.350
\quad\wedge\quad
\mathrm{BAcc}_{\text{no-state}} \le 0.350,
\label{eq:gate}
\end{equation}
and an environment-age row passes only when all three seeds pass.  Chance
balanced accuracy is $0.25$ for the four-way cue.

\section{Method}

\begin{figure}[t]
\centering
\includegraphics[width=\linewidth]{fig_c_architecture.pdf}
\caption{Masked-evidence slot-memory interface and its two data paths.  A causal
render stream is tokenized and written into eight memory slots $M_t$
(stages~1--3).  A label-free target miner withholds salient patch sets and
encodes them with a fixed encoder $\phi$ (stage~4); the slot writer predicts
those hidden targets through a set-prediction loss (stage~5, yellow gradient
path).  Cue labels are never used in training; a post-hoc ridge readout audits
$M_t$ under matched full / reset / no-state conditions (stage~6, dashed
eval-only path).}
\label{fig:architecture}
\end{figure}

Figure~\ref{fig:architecture} is the compact method contract.  A trajectory
cache fixes frames, actions, time indices, and train/validation splits for every
method so that only the memory writer changes.  We describe each component and
then give training and audit pseudocode.

\subsection{Inputs and tokenization}
For each episode the model receives the legal stream: rendered frames
$x_{0:t}$, actions $a_{0:t-1}$, and integer time indices.  A frame CNN maps each
frame to a token $z_i\in\mathbb{R}^{D}$ with $D=160$; learned action and time
embeddings are added, giving a token sequence $Z=[z_0,\dots,z_t]\in
\mathbb{R}^{t\times D}$.  The cue frames are inside the mined-and-hidden set at
training time and outside the length-$K$ window at audit time, so the model
never has a trivially legal copy of the answer.

\subsection{Masked patch-set target}
The training signal is old, visual, and label-free.  For each episode the target
miner splits the causal history $[0,t]$ into $V{=}3$ temporal bins.  Within each
bin it scores patch saliency (local color contrast and novelty against a running
mean), samples one salient view per bin proportionally to saliency, and extracts
a set of high-saliency patches.  The selected patches are \emph{removed from the
visible stream} and encoded by a fixed target encoder $\phi$ into pooled
color/texture latents,
\begin{equation}
Y \;=\; \{\,y_j\,\}_{j=1}^{N}, \qquad
y_j \;=\; \mathrm{stopgrad}\!\big(\phi(\text{patch}_j)\big) \in \mathbb{R}^{D},
\label{eq:target}
\end{equation}
where $\phi$ applies a fixed $4{\times}4$ adaptive pool, mean-centering, and a
saturation channel followed by a frozen random projection.  Because $\phi$ is
fixed and generic, the miner cannot smuggle a semantic label into the target: it
only sees color and texture structure.  Crucially, the target is a \emph{set},
which keeps object-level, spatially distributed structure intact rather than
collapsing it into a single frame reconstruction.

\subsection{Slot writer}
The writer maintains $S{=}8$ normalized memory slots.  Slots act as queries in a
compact multi-head cross-attention over the token sequence, followed by
LayerNorm and an MLP update; stacking these blocks yields the causal memory
state $M_t=[m_1,\dots,m_S]\in\mathbb{R}^{S\times D}$.  Writing one block is
\begin{align}
\tilde{M} &= M + \mathrm{MHA}\big(Q{=}W_q M,\; K{=}W_k Z,\; V{=}W_v Z\big),\\
M' &= \tilde{M} + \mathrm{MLP}\big(\mathrm{LN}(\tilde{M})\big),
\label{eq:slot}
\end{align}
with $H{=}4$ heads.

\subsection{Bounded streaming memory}
\label{sec:streaming}
A causal transformer over the full prefix is causal but not bounded: it can
re-read every historical token at readout.  To make the memory a genuine
bounded persistent state we use a chunked recurrence.  The legal stream is
ingested in consecutive chunks of at most $K$ tokens, and the slots attend only
over the previous slot state concatenated with the current chunk; the chunk
tokens are then \emph{evicted}:
\begin{equation}
M_t \;=\; f_\theta\big(M_{t-1},\; Z_{t-K+1:t}\big),
\qquad
\text{keys/values} = [\,M_{t-1}\,;\,Z_{t-K+1:t}\,],
\label{eq:streaming}
\end{equation}
so all tokens before $t-K+1$ are physically deleted and inaccessible.  Peak
activation memory and per-step compute are therefore constant in sequence
length (Section~\ref{sec:streaming-results}); with $K$ larger than the sequence
the update reduces exactly to the one-shot writer in Eq.~\eqref{eq:slot}, which
we keep as an ablation.  Because both the memory and the target are sets, the
model is not forced to squeeze every old fact into one recurrent vector or one
hand-indexed slot; different slots can specialize to different mined views.

\subsection{Set-prediction objective}
A prediction head $p$ maps slots to predicted target latents
$\hat{Y}=p(M_t)$.  Predictions are matched to the mined target set $Y$ with a
set loss that combines contrastive, cosine, and anti-collapse terms:
\begin{equation}
\mathcal{L}
= \underbrace{\mathcal{L}_{\text{InfoNCE}}(\hat{Y},Y)}_{\text{retain the right facts}}
+\; \lambda_{\cos}\, \underbrace{\big(1-\overline{\cos}(\hat{Y},Y)\big)}_{\text{align latents}}
+\; \lambda_{\text{std}}\, \underbrace{\mathrm{std\text{-}reg}(\hat{Y})}_{\text{anti-collapse}} .
\label{eq:loss}
\end{equation}
The InfoNCE term uses in-batch negatives so the memory must keep the correct old
view separable from other episodes' views; the cosine term aligns matched
latents; the standard-deviation regularizer prevents slot collapse.  No term
uses the cue label.

\subsection{Baselines under a shared contract}
GRU, LSTM, and a plain PyTorch Mamba-lite recurrence replace only the slot
writer.  They consume the identical token sequence and emit a state that is
passed through the same prediction head, same target miner, same encoder $\phi$,
same loss~\eqref{eq:loss}, and same readout.  This is the key experimental
control: any difference in exposure is attributable to the memory format, not to
the data, target, or optimizer.

\subsection{Post-hoc audit readout}
After training, a ridge classifier is fit on \emph{train} memory features and
cue labels, then applied unchanged to held-out full, reset, and no-state
features.  Training is thus deliberately asymmetric: it only applies
reconstruction pressure on hidden patch sets, while evaluation asks a semantic
question that was never optimized directly.  If labels entered training, a
reviewer could attribute success to supervised cue classification; if labels
were never used at all, there would be no way to audit whether the memory
retained the task-relevant cue.  The split between label-free training and
label-revealed evaluation is a measurement design, not a loophole.

\begin{algorithm}[t]
\SetAlgoLined
\DontPrintSemicolon
\KwIn{legal stream $(x_{0:t},a_{0:t-1})$; fixed encoder $\phi$; slot writer $f_\theta$; head $p_\theta$}
\For{each training step}{
  $Z \leftarrow \mathrm{Tokenize}(x_{0:t}, a_{0:t-1}, \tau)$\;
  bins $\leftarrow \mathrm{SplitTemporal}(0{:}t, V{=}3)$\;
  \ForEach{bin $b$}{
    $v_b \leftarrow \mathrm{SampleSalientView}(b)$; \quad patches$_b \leftarrow \mathrm{TopSalientPatches}(v_b)$\;
  }
  $Y \leftarrow \{\,\mathrm{stopgrad}(\phi(\text{patch}))\,\}$;\quad remove mined views from visible $Z$\;
  $M_t \leftarrow f_\theta(Z)$ \tcp*{Eq.~\eqref{eq:slot}, $S{=}8$ slots}
  $\hat{Y} \leftarrow p_\theta(M_t)$; \quad $\mathcal{L} \leftarrow$ Eq.~\eqref{eq:loss} \tcp*{no cue labels}
  update $\theta$ with AdamW on $\nabla_\theta \mathcal{L}$\;
}
\caption{Label-free masked patch-set training (per method; only $f_\theta$ changes).}
\label{alg:train}
\end{algorithm}

\begin{algorithm}[t]
\SetAlgoLined
\DontPrintSemicolon
\KwIn{trained writer $f_\theta$; held-out episodes; cue labels $y$ (eval only)}
$g \leftarrow \mathrm{RidgeFit}(\{f_\theta(Z^{\text{train}})\}, y^{\text{train}})$\;
\ForEach{condition $\in\{$full, reset, no-state$\}$}{
  $M^{\text{cond}} \leftarrow$ memory under condition (clear at $K$ / drop state)\;
  $\mathrm{BAcc}_{\text{cond}} \leftarrow \mathrm{BalancedAcc}(g(M^{\text{cond}}), y^{\text{val}})$\;
}
cell passes $\Leftarrow$ Eq.~\eqref{eq:gate}\;
\caption{Fail-closed audit (labels used only here).}
\label{alg:audit}
\end{algorithm}

\section{Experiments}

\paragraph{Protocol.}
We evaluate twelve rendered OGBench decks \citep{r_ogbench}: PointMaze medium,
large, giant, and teleport; AntMaze large and giant; HumanoidMaze large; Cube
single, double, and triple; Puzzle-3x3; and Scene.  Every memory writer is run
at ages $4/8/15$ and seeds $0/1/2$ ($108$ cells per method).  Figures and tables
are generated from local JSON/NPZ artifacts by \path{scripts/plot_paper_c.py};
per-environment numbers appear in Appendix~\ref{app:full}.  Balanced accuracy is
reported throughout; chance is $0.25$.

\subsection{Architecture comparison under a shared contract}
Table~\ref{tab:main-comparison} and Figure~\ref{fig:baseline-comparison} give
the headline comparison.  The pattern, not any single number, is the claim: full
memory is high while reset and no-state stay near chance, and the recurrent
baselines that solve many short or easy rows lose the exposure pattern on older
or more structured evidence.  Controls are matched at the feature level, so a
deck that leaked the cue through endpoint pixels, actions, class imbalance, or a
forgiving readout would raise the controls and fail.  The slot writer passes all
$36$ environment-age rows and all $108$ seed cells; the strongest recurrent
baseline (GRU) passes $22$ rows and $73$ cells, and LSTM collapses to $6$ rows.

\begin{table}[t]
\centering
\footnotesize
\caption{Controlled-age retention and exposure over 12 decks, 3 ages, and 3
seeds.  Full, reset, and no-state are balanced accuracies from the shared
post-hoc readout; controls are matched on the same held-out rows.}
\label{tab:main-comparison}
\begin{tabular}{lccccc}
\toprule
Method & Full & Reset & No-state & Seed pass & Env-age pass\\
\midrule
\rowcolor{MemCream}Slot memory (ours) & 0.983 & 0.250 & 0.248 & 108/108 & 36/36\\
GRU & 0.803 & 0.252 & 0.250 & 73/108 & 22/36\\
LSTM & 0.577 & 0.248 & 0.250 & 26/108 & 6/36\\
Mamba-lite & 0.800 & 0.255 & 0.253 & 71/108 & 22/36\\
\bottomrule
\end{tabular}

\end{table}

\begin{figure}[t]
\centering
\includegraphics[width=0.86\linewidth]{fig_c_baseline_comparison.pdf}
\caption{Architecture comparison with all non-memory components held fixed.  The
slot writer dominates mean retention, the age-15 stress test, and the
fail-closed all-pass count.}
\label{fig:baseline-comparison}
\end{figure}

\subsection{Shortcut diagnostics}
Figure~\ref{fig:cue-inference} makes the control behavior explicit on the
hardest deck (Puzzle age~15).  In the cue-inference panels (a), full memory
concentrates readout evidence on the true old cue, while reset and no-state are
not cleanly random: they collapse toward endpoint-biased columns, which is
exactly the shortcut the audit is designed to expose.  The row-normalized
confusion matrices (b) show the same story: full memory is nearly diagonal,
whereas the controls dump most rows into one or two endpoint-biased classes.
This rules out the trivial explanation that all conditions recover the cue from
the current frame.

\begin{figure}[t]
\centering
\includegraphics[width=\linewidth]{fig_c_cue_inference.pdf}
\caption{Shortcut diagnostics on held-out Puzzle age-15 features.
(a)~Cue inference: dashed boxes mark the true cue, solid boxes the prediction;
full memory tracks the old cue while the controls collapse toward endpoint bias.
(b)~Row-normalized confusion: full memory is diagonal, the controls are not.}
\label{fig:cue-inference}
\end{figure}

\subsection{Evidence-age stress and unseen-age generalization}
The age sweep is a stress test (Table~\ref{tab:age-comparison},
Figure~\ref{fig:age-gen}a).  Age~4 asks whether evidence just outside the
context boundary is accessible; age~8 tests whether memory is more than a short
transient; age~15 is the long-delay probe.  The useful signal is the curve
shape.  LSTM decays first; GRU and Mamba-lite stay stronger but fail many
long-age rows; the slot writer holds a stable exposure margin.  Age~15 is
therefore the architectural separator, not a cosmetic extra setting.

We further ask whether the interface is tied to the specific registered ages.
We train the slot writer on a \emph{mixture} of ages $\{4,8,15\}$ and evaluate
on a denser grid $\{4,6,8,10,12,15,18\}$, so that $\{6,10,12,18\}$ are never
seen during training.  Table~\ref{tab:auto-age} and Figure~\ref{fig:age-gen}b
show that every held-out age clears the same fail-closed gate, with full-memory
exposure indistinguishable from the trained ages and both controls at chance.
The delay at which evidence must be retained is therefore not memorized as a
fixed offset; mixed-age training yields a delay-robust interface.  We are
careful about the boundary: this demonstrates unseen-age \emph{generalization},
not a learned explicit age router or an automatic maximum-age search, which we
leave to future work (Appendix~\ref{app:autoage}).

\begin{table}[t]
\centering
\footnotesize
\begin{minipage}[t]{0.46\linewidth}
\centering
\caption{Age-wise mean full-memory balanced accuracy (fixed-age training).}
\label{tab:age-comparison}
\input{generated_results/age_comparison_tabular.tex}
\end{minipage}\hfill
\begin{minipage}[t]{0.50\linewidth}
\centering
\caption{Mixed-age training (ages $\{4,8,15\}$), evaluated on a denser grid;
held-out ages shaded.}
\label{tab:auto-age}
\resizebox{\linewidth}{!}{\begin{tabular}{lccccc}
\toprule
Eval age & Trained? & Pass & Full & Reset & No-state\\
\midrule
4 & yes & 36/36 & 0.984 & 0.246 & 0.253\\
6 & held-out & 36/36 & 0.985 & 0.254 & 0.247\\
8 & yes & 36/36 & 0.985 & 0.252 & 0.250\\
10 & held-out & 36/36 & 0.988 & 0.253 & 0.247\\
12 & held-out & 36/36 & 0.988 & 0.255 & 0.255\\
15 & yes & 36/36 & 0.985 & 0.250 & 0.249\\
18 & held-out & 36/36 & 0.982 & 0.248 & 0.251\\
\bottomrule
\end{tabular}
}
\end{minipage}
\end{table}

\begin{figure}[t]
\centering
\includegraphics[width=\linewidth]{fig_c_age_generalization.pdf}
\caption{(a)~Fixed-age training: recurrent baselines decay as evidence ages,
while the slot writer stays above the gate.  (b)~Mixed-age training generalizes
to unseen ages (diamonds); full-memory exposure stays high while reset and
no-state controls remain at chance.}
\label{fig:age-gen}
\end{figure}

\subsection{Where the gap concentrates}
Figure~\ref{fig:env-heatmaps} replaces a cherry-picked hard-row table with a
full visual analysis over all decks.  Panel~(a) shows that some locomotion-like
rows have small gaps, which argues against an unfairly weak comparison: the
baselines are competent when the evidence is short and coarse.  The larger gaps
appear in Puzzle, Scene, Cube, and long-age PointMaze rows, where object
identity, order, or spatial configuration must stay separated until readout.
Panel~(b) confirms that slot success is not carried by a single outlier: the
audit margin stays well above chance across the registered suite, with the
hardest pressure on Humanoid, Puzzle, and Scene at age~15.

\begin{figure}[t]
\centering
\includegraphics[width=\linewidth]{fig_c_env_heatmaps.pdf}
\caption{Per-environment analysis over all twelve decks and three ages.
(a)~Slot memory minus the best recurrent baseline (full BAcc gap): gains
concentrate on structured and long-age rows.  (b)~Slot audit margin, full
$-\max(\text{reset},\text{no-state})$: every registered cell clears the gate.}
\label{fig:env-heatmaps}
\end{figure}

\subsection{Bounded native use}
The native-use stage asks a different question from the readout audit
(Figure~\ref{fig:native-use}, Table~\ref{tab:native-use}).  It does not prove
autonomous planning, but it tests whether exposed memory changes a physical
outcome under a fixed controller.  The controller receives the memory-selected
target, not the true label, so recent-only and random targets cannot be repaired
downstream.  Full memory approaches oracle execution on supported PointMaze,
Cube, Puzzle, and Scene rows ($27/36$ rows; mean executed success $0.957$; worst
scored row $0.828$), a mean lift of about $0.68$ over both the recent-only and
random-target controls.  AntMaze and HumanoidMaze are outside the native-use
claim because the repository does not yet contain audited local low-level
controllers for those bodies; those rows are reported as controller-unavailable
in Appendix~\ref{app:full}, not as failures.  We treat this stage as a
\emph{bounded-use sanity check}: it shows that a correctly exposed memory label
selects a target a fixed controller can execute, not that the memory performs
world-model planning.  The stronger world-model evidence is the frozen-LeWM
integration in Section~\ref{sec:revision}, where the memory changes the host's
own predictive output.

\begin{figure}[t]
\centering
\includegraphics[width=0.78\linewidth]{fig_c_native_heatmap.pdf}
\caption{Fixed-controller native use on supported decks.  Each cell shows
full-memory executed success and, below it, the gain over the stronger
recent-only or random-target control.}
\label{fig:native-use}
\end{figure}

\begin{table}[t]
\centering
\footnotesize
\caption{Fixed-controller native-use summary for supported OGBench rows.}
\label{tab:native-use}
\begin{tabular}{lc}
\toprule
Metric & Value\\
\midrule
Scored env-age rows & 27/36\\
Mean full-memory executed success & 0.957\\
Worst scored full-memory row & 0.828\\
Mean recent-only success & 0.279\\
Mean random-target success & 0.279\\
Mean oracle success & 0.973\\
\bottomrule
\end{tabular}

\end{table}

\section{Bounded memory, capacity, robustness, and world-model integration}
\label{sec:revision}

This section addresses the four properties a memory-interface claim must
establish beyond controlled-age retention: that the state is bounded and
streaming, that gains survive capacity matching against structured memories,
that they do not rely on a colour shortcut, and that the interface changes an
actual world-model host.

\subsection{The memory is bounded and streaming}
\label{sec:streaming-results}
We reformulate the writer as the chunked recurrence of Eq.~\eqref{eq:streaming}
and evict each chunk after it is consumed.  Figure~\ref{fig:streaming-cost}
measures the writer as the observed sequence grows.  The one-shot full-prefix
writer's peak activation memory grows with length, while the streaming writer's
per-step compute is flat and its peak memory stays bounded: the persistent
state is a fixed $S\times D$ tensor and no pre-eviction tokens are retained.
The audit numbers in the previous section are reported with this streaming
writer, so the retention claim is a property of a bounded state rather than of
re-reading history.

\begin{figure}[t]
\centering
\includegraphics[width=0.86\linewidth]{fig_c_streaming_cost.pdf}
\caption{Bounded streaming memory.  (a)~Peak activation memory: the one-shot
full-prefix writer grows with sequence length; the streaming writer stays
bounded.  (b)~Per-step compute is constant for the streaming writer.}
\label{fig:streaming-cost}
\end{figure}

\subsection{Gains survive capacity matching and structured baselines}
Because the audit readout consumes a mean-pooled $D$-dimensional vector for
\emph{every} method, the comparison is already dimension-matched at the
readout.  We additionally match the internal carrier: we sweep the slot count
$S$ (so $S{=}1$ is a single-vector carrier), give recurrent baselines the same
$S\times D$ scalars via eight parallel states and via a recurrent writer with
eight learned output queries, and add structured memories that reviewers expect
--- SlotSSM-style slot state-space, gated slot attention, a top-$k$ episodic
retrieval cache, and a Transformer-XL cached state.
Figure~\ref{fig:capacity} plots age-15 full-memory accuracy against exposed
state scalars and compares all carriers under the shared contract.  The slot
interface remains at the top of the accuracy-vs-capacity frontier: matching the
carrier size does not close the gap, and structured recurrent/retrieval
memories do not match the slot writer's long-age exposure at equal state size.

\begin{figure}[t]
\centering
\includegraphics[width=\linewidth]{fig_c_capacity_structured.pdf}
\caption{Capacity-matched and structured comparison at age~15.  (a)~Full-memory
balanced accuracy across recurrent, state-space, slot, retrieval, and cached
baselines under the identical target/loss/readout contract.  (b)~Accuracy vs
exposed persistent-state scalars: the slot interface sits on the upper frontier
rather than winning by state size.}
\label{fig:capacity}
\end{figure}

\begin{table}[t]
\centering
\footnotesize
\caption{Structured and capacity-matched memory baselines (age~15, matched
$D{=}160$).  State scalars are the exposed persistent-state size.}
\label{tab:structured}
\begin{tabular}{lcccc}
\toprule
Memory & Full & Reset & No-state & State scalars\\
\midrule
\rowcolor{MemCream}Slot (ours) & 0.735 & 0.248 & 0.260 & 1280\\
SlotSSM & 0.451 & 0.266 & 0.244 & 1280\\
GSA & 0.669 & 0.214 & 0.267 & 1280\\
Top-k retr. & 0.657 & 0.250 & 0.232 & 1280\\
Transf.-XL & 0.294 & 0.293 & 0.308 & 1920\\
8x GRU & 0.326 & 0.251 & 0.258 & 1280\\
GRU+queries & 0.626 & 0.244 & 0.226 & 160\\
GRU & 0.400 & 0.260 & 0.259 & 160\\
LSTM & 0.293 & 0.275 & 0.274 & 320\\
Mamba-lite & 0.515 & 0.282 & 0.269 & 160\\
\bottomrule
\end{tabular}

\end{table}

\subsection{Retention is not a colour shortcut}
The default cue makes the label recoverable from pooled colour, which the fixed
target encoder preserves.  We therefore add a shortcut-broken benchmark in
which the label is the cue \emph{shape identity} while its colour is randomized
independently of the label, so pooled colour is uninformative.
Figure~\ref{fig:shortcut} shows that the slot interface still exposes the old
identity well above the matched controls, while the controls stay at chance:
the memory retains object-like structure, not merely a hue.

\begin{figure}[t]
\centering
\includegraphics[width=0.52\linewidth]{fig_c_shortcut.pdf}
\caption{Randomized-colour identity cue: colour is decorrelated from the label,
so pooled colour cannot solve the task.  Full memory still separates the old
cue while reset and no-state remain at chance.}
\label{fig:shortcut}
\end{figure}

\subsection{Integration into a frozen LeWM world-model host}
To show the interface changes an actual world model rather than a standalone
probe, we inject the slot memory as a carrier into the frozen official LeWM
predictor \citep{r_lewm} on a PushT visual-binding task, where the decisive cue
lies outside the host's context window.  We measure whether the frozen host's
\emph{own output} exposes the old evidence under matched full/reset/no-state
controls, with cue labels used only for the post-hoc readout.
Figure~\ref{fig:lewm} shows that the injected memory raises host-output balanced
accuracy far above both controls on a six-way binding task at age~15 across five
seeds, clearing the same fail-closed gate.  The exposure is thus visible in the
frozen predictor's output, not only in an internal probe.

\begin{figure}[t]
\centering
\includegraphics[width=0.42\linewidth]{fig_c_lewm_host.pdf}
\caption{Frozen-LeWM integration.  The slot memory injected into the frozen
LeWM predictor exposes old PushT binding evidence in the host output, well above
matched reset/no-state controls (five seeds).}
\label{fig:lewm}
\end{figure}

\subsection{Age scaling}
Finally we test how far back the interface retains evidence.  We re-render
longer PointMaze-large episodes, train on ages $\{4,8,16\}$, and evaluate on a
geometric ladder up to age~128 at fixed per-step budget
(Figure~\ref{fig:age-scaling}).  Full-memory exposure decays gracefully with a
long empirical half-life while the controls stay at chance, quantifying the
maximum passing age rather than asserting a single stress point.

\begin{figure}[t]
\centering
\includegraphics[width=0.5\linewidth]{fig_c_age_scaling.pdf}
\caption{Age scaling to 128 on re-rendered PointMaze-large: trained ages
(circles) and unseen ages (diamonds) at fixed per-step memory/compute.}
\label{fig:age-scaling}
\end{figure}

\section{Discussion}

The experimental picture is coherent.  The target miner creates an old,
label-free visual prediction problem; the slot writer provides a memory format
that matches the target's multi-part structure; the readout audit shows that old
evidence is exposed rather than hidden inside an inaccessible state; and the
fixed-controller stage shows that exposed evidence can alter execution when a
valid controller exists.  Three insights follow.

\paragraph{Structure-preserving targets matter more than another recurrent block.}
Flat whole-frame reconstruction can look like memory while remaining too diffuse
for object-level cues; once the target becomes a \emph{set} of salient patches,
the same slot bottleneck becomes trainable without cue labels.  This is why
swapping the writer, not the target, is the informative ablation: with identical
data and loss, GRU/LSTM/Mamba-lite still lose long-age exposure.

\paragraph{Short-horizon competence is not long-horizon exposure.}
GRU and Mamba-lite often clear easy age-4 rows, yet the age-15 curve and the
Puzzle/Scene/Cube gaps show where a single hidden vector stops keeping separable
visual facts.  Reporting only mean accuracy would hide this; the age sweep and
the per-environment gap heatmap are what localize the failure.

\paragraph{Internal storage is not host exposure.}
The fail-closed controls enforce a distinction that a raw memory-capacity metric
misses: a carrier can score well on an internal probe while a frozen interface
still fails to expose old evidence as evidence age grows.  Our audit therefore
reports retention, exposure, and use as separate rungs rather than a single
pooled score.

\paragraph{Delay robustness is a property of the interface, not the offset.}
Mixed-age training generalizes to unseen delays
(Table~\ref{tab:auto-age}, Figure~\ref{fig:age-gen}b), so the exposure is not a
memorized offset.  The boundary is still explicit: this is unseen-age
generalization, not a learned age router that decides \emph{how far back} to
retrieve, nor an automatic maximum-age search; those remain future work
(Appendix~\ref{app:autoage}).  The clean claim is a controlled memory-interface
result: a masked-evidence slot interface retains and exposes old visual evidence
more reliably than recurrent alternatives trained with the same objective, it
generalizes across evidence delays, and the exposed state can drive
fixed-controller execution where controllers exist.

\section{Conclusion}

Old-evidence memory should be evaluated by retention, exposure, and bounded use.
A hidden state that stores information but never exposes it is not enough for a
frozen world-model interface; a readout that succeeds without controls is not
enough either.  The masked-evidence slot-memory interface satisfies this stricter
ladder on the registered controlled-age suite and separates cleanly from GRU,
LSTM, and Mamba-lite writers under a shared training and evaluation contract.
The result supports an ICLR-style controlled-age memory paper when the claim
remains scoped to structured retention and exposure.

\section{Limitations}

The benchmark is controlled by design.  Cues are synthetic overlays on rendered
OGBench streams, which enables shortcut audits but does not cover all real robot
memory patterns.  The native-use claim depends on fixed controllers and is
reported only for environments with audited local controllers.  The Mamba-lite
baseline tests selective recurrence but is not the optimized official Mamba
stack.  The readout uses labels after training, so this is an audit of exposed
representation quality rather than an end-to-end policy.  Most importantly,
automatic age discovery remains future work: the current paper tests registered
delays, while the next stage must learn a delay-robust router and report
unseen-age generalization.
